\documentclass[12pt,a4paper]{article}
\usepackage[utf8]{inputenc}
\usepackage[english]{babel}
\usepackage{amsmath}
\usepackage{amsfonts}
\usepackage{amssymb}
\usepackage{graphicx}
\usepackage{hyperref}
\usepackage{algorithm}
\usepackage{algorithmic}
\usepackage{listings}
\usepackage{xcolor}
\usepackage{geometry}
\usepackage{booktabs}
\usepackage{subfig}
\usepackage{cite}

\geometry{margin=1in}

\title{Sensor Fusion for Navigation Systems:\\
Implementation and Theory}
\author{Sensor Fusion Project}
\date{\today}

\lstset{
    language=C++,
    basicstyle=\ttfamily\small,
    keywordstyle=\color{blue},
    commentstyle=\color{gray},
    stringstyle=\color{red},
    numbers=left,
    numberstyle=\tiny,
    frame=single,
    breaklines=true
}

\begin{document}

\maketitle
\tableofcontents
\newpage

\section{Introduction}

This document provides comprehensive documentation for a sensor fusion navigation system that integrates data from multiple sensors including Inertial Measurement Units (IMU), Global Positioning System (GPS), barometric pressure sensors, and magnetometers. The implementation is designed for cross-platform deployment, supporting both x64 simulation environments and embedded systems.

\subsection{Project Goals}

\begin{itemize}
    \item Implement accurate sensor models with realistic error characteristics
    \item Develop Extended Kalman Filter (EKF) for optimal state estimation
    \item Support cross-compilation for embedded platforms
    \item Provide comprehensive documentation based on established literature
    \item Create modular, testable, and maintainable C++ code
\end{itemize}

\subsection{System Overview}

The sensor fusion system combines complementary sensor measurements to estimate the navigation state (position, velocity, and attitude) of a moving platform. Each sensor provides different information with varying accuracy and update rates:

\begin{itemize}
    \item \textbf{IMU}: High-rate (100-1000 Hz) acceleration and angular velocity measurements
    \item \textbf{GPS}: Low-rate (1-10 Hz) absolute position and velocity with bounded errors
    \item \textbf{Barometer}: Medium-rate (10-50 Hz) altitude measurements
    \item \textbf{Magnetometer}: Medium-rate (10-100 Hz) heading reference
\end{itemize}

\section{Theoretical Background}

\subsection{Navigation State Representation}

The navigation state vector consists of position, velocity, and attitude:

\begin{equation}
\mathbf{x} = \begin{bmatrix}
\mathbf{p}^n \\
\mathbf{v}^n \\
\mathbf{q} \\
\mathbf{b}_a \\
\mathbf{b}_g
\end{bmatrix}
\end{equation}

where:
\begin{itemize}
    \item $\mathbf{p}^n = [p_N, p_E, p_D]^T$ is position in North-East-Down (NED) frame
    \item $\mathbf{v}^n = [v_N, v_E, v_D]^T$ is velocity in NED frame
    \item $\mathbf{q} = [q_w, q_x, q_y, q_z]^T$ is orientation quaternion
    \item $\mathbf{b}_a = [b_{ax}, b_{ay}, b_{az}]^T$ is accelerometer bias
    \item $\mathbf{b}_g = [b_{gx}, b_{gy}, b_{gz}]^T$ is gyroscope bias
\end{itemize}

\subsection{Inertial Measurement Unit (IMU)}

\subsubsection{IMU Measurement Model}

The IMU provides measurements of specific force and angular velocity in the body frame:

\begin{equation}
\tilde{\mathbf{f}} = \mathbf{f} + \mathbf{b}_a + \mathbf{n}_a
\end{equation}

\begin{equation}
\tilde{\boldsymbol{\omega}} = \boldsymbol{\omega} + \mathbf{b}_g + \mathbf{n}_g
\end{equation}

where $\mathbf{n}_a$ and $\mathbf{n}_g$ are white Gaussian noise processes.

\subsubsection{Error Models}

MEMS IMU errors are typically classified as \cite{aggarwal2013comparison}:

\begin{itemize}
    \item \textbf{Constant bias}: Fixed offset, determined during calibration
    \item \textbf{Bias instability}: Slow-varying bias modeled as random walk or first-order Gauss-Markov process
    \item \textbf{White noise}: High-frequency random noise
    \item \textbf{Scale factor error}: Multiplicative error in sensitivity
    \item \textbf{Cross-axis sensitivity}: Non-orthogonality of sensor axes
    \item \textbf{Temperature effects}: Temperature-dependent bias and scale factor
\end{itemize}

The bias evolution is modeled as:

\begin{equation}
\dot{\mathbf{b}}_a = -\frac{1}{\tau_a}\mathbf{b}_a + \mathbf{w}_a
\end{equation}

\begin{equation}
\dot{\mathbf{b}}_g = -\frac{1}{\tau_g}\mathbf{b}_g + \mathbf{w}_g
\end{equation}

where $\tau$ is the correlation time and $\mathbf{w}$ is white noise driving the random walk.

\subsection{Strapdown Navigation Equations}

The mechanization equations for strapdown inertial navigation are \cite{savage1998strapdown, titterton2004strapdown}:

\begin{equation}
\dot{\mathbf{p}}^n = \mathbf{v}^n
\end{equation}

\begin{equation}
\dot{\mathbf{v}}^n = \mathbf{C}_b^n \mathbf{f}^b + \mathbf{g}^n - (2\boldsymbol{\omega}_{ie}^n + \boldsymbol{\omega}_{en}^n) \times \mathbf{v}^n
\end{equation}

\begin{equation}
\dot{\mathbf{q}} = \frac{1}{2}\mathbf{q} \otimes \boldsymbol{\omega}_{nb}^b
\end{equation}

where:
\begin{itemize}
    \item $\mathbf{C}_b^n$ is the direction cosine matrix from body to navigation frame
    \item $\mathbf{g}^n$ is gravity in the navigation frame
    \item $\boldsymbol{\omega}_{ie}^n$ is Earth rotation rate
    \item $\boldsymbol{\omega}_{en}^n$ is transport rate
    \item $\otimes$ denotes quaternion multiplication
\end{itemize}

\subsection{GPS Measurement Model}

GPS provides measurements of position and velocity with errors:

\begin{equation}
\mathbf{z}_{GPS,p} = \mathbf{p}^n + \mathbf{n}_{GPS,p}
\end{equation}

\begin{equation}
\mathbf{z}_{GPS,v} = \mathbf{v}^n + \mathbf{n}_{GPS,v}
\end{equation}

Typical GPS accuracy (95\% confidence):
\begin{itemize}
    \item Horizontal position: 2-5 meters
    \item Vertical position: 3-8 meters
    \item Velocity: 0.1-0.3 m/s
\end{itemize}

\subsection{Barometric Altimeter}

The barometric altimeter measures altitude based on atmospheric pressure:

\begin{equation}
h = 44330 \left(1 - \left(\frac{P}{P_0}\right)^{0.1903}\right)
\end{equation}

where $P$ is measured pressure and $P_0$ is sea-level pressure (101325 Pa).

Measurement model:
\begin{equation}
\mathbf{z}_{baro} = -p_D + \mathbf{n}_{baro}
\end{equation}

Typical accuracy: 1-3 meters in stable conditions.

\subsection{Magnetometer}

The magnetometer measures the local magnetic field:

\begin{equation}
\tilde{\mathbf{m}}^b = \mathbf{C}_n^b \mathbf{m}^n + \mathbf{b}_m + \mathbf{n}_m
\end{equation}

where $\mathbf{m}^n$ is the local magnetic field in the navigation frame.

\section{Extended Kalman Filter}

\subsection{Filter Formulation}

The EKF provides optimal state estimation for the nonlinear navigation equations \cite{bar2001estimation}. The filter operates in two phases:

\subsubsection{Prediction (Time Update)}

\begin{equation}
\hat{\mathbf{x}}_{k|k-1} = f(\hat{\mathbf{x}}_{k-1|k-1}, \mathbf{u}_k)
\end{equation}

\begin{equation}
\mathbf{P}_{k|k-1} = \mathbf{F}_k \mathbf{P}_{k-1|k-1} \mathbf{F}_k^T + \mathbf{Q}_k
\end{equation}

where:
\begin{itemize}
    \item $\mathbf{F}_k$ is the state transition matrix (Jacobian of $f$)
    \item $\mathbf{Q}_k$ is the process noise covariance
\end{itemize}

\subsubsection{Update (Measurement Update)}

\begin{equation}
\mathbf{K}_k = \mathbf{P}_{k|k-1} \mathbf{H}_k^T (\mathbf{H}_k \mathbf{P}_{k|k-1} \mathbf{H}_k^T + \mathbf{R}_k)^{-1}
\end{equation}

\begin{equation}
\hat{\mathbf{x}}_{k|k} = \hat{\mathbf{x}}_{k|k-1} + \mathbf{K}_k (\mathbf{z}_k - h(\hat{\mathbf{x}}_{k|k-1}))
\end{equation}

\begin{equation}
\mathbf{P}_{k|k} = (\mathbf{I} - \mathbf{K}_k \mathbf{H}_k) \mathbf{P}_{k|k-1}
\end{equation}

where:
\begin{itemize}
    \item $\mathbf{H}_k$ is the measurement matrix (Jacobian of $h$)
    \item $\mathbf{R}_k$ is the measurement noise covariance
    \item $\mathbf{K}_k$ is the Kalman gain
\end{itemize}

\subsection{State Transition Matrix}

For the INS error-state model, the state transition matrix is:

\begin{equation}
\mathbf{F} = \begin{bmatrix}
\mathbf{0} & \mathbf{I} & \mathbf{0} & \mathbf{0} & \mathbf{0} \\
\mathbf{0} & -2[\boldsymbol{\omega}_{ie}^n]_\times & [\mathbf{C}_b^n\mathbf{f}^b]_\times & \mathbf{C}_b^n & \mathbf{0} \\
\mathbf{0} & \mathbf{0} & -[\boldsymbol{\omega}_{ib}^b]_\times & \mathbf{0} & -\mathbf{C}_b^n \\
\mathbf{0} & \mathbf{0} & \mathbf{0} & -\frac{1}{\tau_a}\mathbf{I} & \mathbf{0} \\
\mathbf{0} & \mathbf{0} & \mathbf{0} & \mathbf{0} & -\frac{1}{\tau_g}\mathbf{I}
\end{bmatrix}
\end{equation}

where $[\cdot]_\times$ denotes the skew-symmetric matrix operator.

\subsection{Measurement Models}

\subsubsection{GPS Position Measurement}

\begin{equation}
\mathbf{H}_{GPS,p} = \begin{bmatrix} \mathbf{I}_3 & \mathbf{0} & \mathbf{0} & \mathbf{0} & \mathbf{0} \end{bmatrix}
\end{equation}

\subsubsection{GPS Velocity Measurement}

\begin{equation}
\mathbf{H}_{GPS,v} = \begin{bmatrix} \mathbf{0} & \mathbf{I}_3 & \mathbf{0} & \mathbf{0} & \mathbf{0} \end{bmatrix}
\end{equation}

\subsubsection{Barometer Measurement}

\begin{equation}
\mathbf{H}_{baro} = \begin{bmatrix} 0 & 0 & -1 & \mathbf{0}_{1\times12} \end{bmatrix}
\end{equation}

\section{Implementation Details}

\subsection{Software Architecture}

The implementation follows object-oriented design principles with clear separation of concerns:

\begin{itemize}
    \item \textbf{Sensor classes}: Encapsulate sensor physics and error models
    \item \textbf{Fusion classes}: Implement Kalman filter algorithms
    \item \textbf{Utility classes}: Provide coordinate transformations and math operations
\end{itemize}

\subsection{Key Design Features}

\begin{enumerate}
    \item \textbf{Header-only matrix library}: Using Eigen3 for efficient linear algebra
    \item \textbf{Modular sensor interface}: Easy to add new sensor types
    \item \textbf{Configurable error models}: Runtime parameter adjustment
    \item \textbf{Cross-compilation support}: CMake toolchain files for ARM targets
    \item \textbf{Minimal dependencies}: Suitable for embedded deployment
\end{enumerate}

\subsection{Build Instructions}

\subsubsection{x64 Build}

\begin{lstlisting}[language=bash]
mkdir build && cd build
cmake ..
make
make test
\end{lstlisting}

\subsubsection{ARM Cross-Compilation}

\begin{lstlisting}[language=bash]
mkdir build-arm && cd build-arm
cmake -DCMAKE_TOOLCHAIN_FILE=../toolchain-arm.cmake ..
make
\end{lstlisting}

\section{Sensor Specifications}

\subsection{IMU Specifications}

Typical MEMS IMU characteristics (consumer-grade):

\begin{table}[h]
\centering
\begin{tabular}{lcc}
\toprule
Parameter & Accelerometer & Gyroscope \\
\midrule
Range & $\pm 16$ g & $\pm 2000$ deg/s \\
Noise density & $150$ $\mu$g/$\sqrt{Hz}$ & $0.01$ deg/s/$\sqrt{Hz}$ \\
Bias instability & $40$ $\mu$g & $10$ deg/h \\
Scale factor error & $0.5$ \% & $0.5$ \% \\
Temperature coefficient & $0.01$ \%/°C & $0.01$ \%/°C \\
\bottomrule
\end{tabular}
\caption{Typical MEMS IMU specifications}
\end{table}

\subsection{GPS Specifications}

\begin{table}[h]
\centering
\begin{tabular}{lc}
\toprule
Parameter & Value \\
\midrule
Horizontal accuracy (95\%) & 2.5 m \\
Vertical accuracy (95\%) & 5.0 m \\
Velocity accuracy (95\%) & 0.1 m/s \\
Update rate & 1-10 Hz \\
Time to first fix & 30 s (cold start) \\
\bottomrule
\end{tabular}
\caption{GPS receiver specifications}
\end{table}

\section{Validation and Testing}

\subsection{Unit Tests}

Individual sensor models and filter components are tested:
\begin{itemize}
    \item IMU error model validation
    \item Coordinate transformation accuracy
    \item EKF numerical stability
    \item Sensor fusion convergence
\end{itemize}

\subsection{Simulation Scenarios}

\begin{enumerate}
    \item \textbf{Stationary test}: Verify bias estimation
    \item \textbf{Constant velocity}: Validate basic integration
    \item \textbf{Turning maneuver}: Test attitude estimation
    \item \textbf{GPS outage}: Evaluate dead-reckoning performance
    \item \textbf{3D trajectory}: Full system validation
\end{enumerate}

\section{Performance Analysis}

\subsection{Error Propagation}

Without GPS updates, INS errors grow according to \cite{woodman2007introduction}:

\begin{itemize}
    \item Position error: $\propto t^3$ (gyro bias) and $\propto t^2$ (accel bias)
    \item Velocity error: $\propto t^2$ (gyro bias) and $\propto t$ (accel bias)
    \item Attitude error: $\propto t$ (gyro bias)
\end{itemize}

\subsection{Fusion Benefits}

Sensor fusion provides:
\begin{itemize}
    \item GPS bounds long-term drift
    \item IMU provides high-rate, smooth state estimates
    \item Barometer improves vertical accuracy
    \item Magnetometer provides heading reference
\end{itemize}

\section{Future Enhancements}

\subsection{Algorithm Improvements}

\begin{itemize}
    \item Unscented Kalman Filter (UKF) implementation \cite{wan2000unscented}
    \item Adaptive covariance tuning
    \item Multi-rate sensor handling
    \item Outlier rejection and fault detection
\end{itemize}

\subsection{Additional Sensors}

\begin{itemize}
    \item Visual odometry integration
    \item Wheel encoder fusion (ground vehicles)
    \item Ultra-wideband (UWB) ranging
    \item LiDAR-based localization
\end{itemize}

\section{Conclusion}

This document describes a comprehensive sensor fusion system for navigation. The implementation balances theoretical rigor with practical considerations for embedded deployment. The modular architecture enables easy extension and customization for specific applications.

\bibliographystyle{ieeetr}
\bibliography{references}

\end{document}
